\section{The completed cohomology calculation}
\label{sec:analytic}

The goal is the completed coefficient calculation in
Theorem~\ref{thm:completed-calculation}.  We first compute the geometric
cohomology and its actions, then descend Frobenius at each finite marking
stage and take the marking colimit.  Throughout this section, \(C\) is an
algebraically closed complete perfectoid extension of \(\mathbb C_p\),
with a chosen embedding \(k\hookrightarrow C^\flat\).  On diamonds over
\(C^\flat\) we write \(\mathcal O^\flat\) for the characteristic-\(p\)
structure sheaf.  Thus, on the diamond of a perfectoid space in
characteristic \(p\), this is its usual structure sheaf; on an untilted
analytic space it is the structure sheaf on its characteristic-\(p\)
perfectoid site.

Let \(T\) be an arbitrary compact profinite set.  Products with \(T\)
always mean products with its associated locally profinite v-sheaf.
For a complete topological coefficient group \(M\), put
\[
 C^r_{\mathrm{cts}}(G,M;T)
   =\operatorname{Map}_{\mathrm{cts}}(T\times G^r,M).
\]
When \(G\) is compact and \(M\) is a Banach space, this mapping space has
the supremum norm.  For a Fr\'echet space it has the supremum seminorms
of a countable family of Banach seminorms.  Products indexed by a
discrete set have their product topology; they carry no boundedness
condition across that discrete set.  These conventions will be
preserved throughout the calculations.

\subsection{Two facts about continuous families and countable limits}

We first record the approximation argument used in the analytic
exhaustions.

\begin{samepage}
\begin{lemma}
\label{lem:analytic-roos}
Let
\[
 M_0 \xleftarrow{f_0} M_1 \xleftarrow{f_1} M_2
       \xleftarrow{} \cdots
\]
be a countable inverse system of complete nonarchimedean normed
abelian groups.  Suppose every \(f_n\) is a contraction with dense
image.  Then the map
\begin{equation}
 \prod_{n\geq0}M_n\longrightarrow\prod_{n\geq0}M_n,\qquad
 (x_n)\longmapsto(x_n-f_n(x_{n+1}))
 \label{eq:analytic-roos}
\end{equation}
is surjective.  The same assertion holds after replacing every \(M_n\)
by \(\operatorname{Map}_{\mathrm{cts}}(T,M_n)\).
Consequently these systems have no positive derived inverse limits.
\end{lemma}
\end{samepage}

\begin{proof}
Fix a right-hand side \((y_n)\) and positive real numbers
\(\varepsilon_n\to0\).  Suppose that a finite tuple
\(x_0,\ldots,x_n\) solves the first \(n\) equations.  By density choose
\(x_{n+1}\) such that
\[
 d_n=y_n+f_n(x_{n+1})-x_n,\qquad \|d_n\|<\varepsilon_n.
\]
Replace \(x_n\) by \(x_n+d_n\), and replace each earlier \(x_i\) by
\[
 x_i+f_i f_{i+1}\cdots f_{n-1}(d_n).
\]
The enlarged tuple now solves the first \(n+1\) equations.  Each
previously chosen coordinate has changed by at most
\(\varepsilon_n\).  Starting with any \(x_0\), this procedure produces
a Cauchy sequence in every fixed coordinate.  Completeness gives a
solution of all equations.

For the parameter assertion, the mapping spaces are complete and
the induced maps are contractions.  They have dense image: a
continuous map from \(T\) to \(M_n\) is uniformly approximated by a map
constant on a finite clopen partition, and its finitely many values
can be approximated in the image of \(M_{n+1}\).  Apply the preceding
argument to these mapping spaces.  Finally, the two-term Roos
complex \eqref{eq:analytic-roos} computes the derived inverse limit
of a countable system.
\end{proof}

The \(c_0\)-condition on a weighted coefficient family means that,
for every \(\varepsilon>0\), only finitely many weighted coefficient
norms are at least \(\varepsilon\).
We use two other elementary consequences of compactness.  First, a
continuous map from \(T\) to a weighted \(c_0\)-space has uniformly
small tails.  Indeed, its compact image admits a finite cover by
small balls, and the finitely many centers have uniformly small
tails.  Conversely, continuous coordinate functions with uniformly
small tails define a continuous map into that \(c_0\)-space.  Applying
this observation to each seminorm gives its Fr\'echet version.
Second, if \(V\) is discrete, every continuous map \(T\to V\) has
finite image.  In particular, continuous \(V\)-valued functions lift
through every surjection of discrete groups, by choosing lifts on
a finite clopen partition.

For use with Kummer covers, continuous cohomology of
\(\mathbb Z_p\) on a complete, separated, linearly topologized
\(\Fp\)-module \(M\) with invariant neighborhoods is computed by
\begin{equation}
 [\,M\xrightarrow{\gamma-1}M\,],
 \label{eq:analytic-cyclic}
\end{equation}
in degrees \(0,1\), where \(\gamma\) is a topological generator.
Here one applies continuous \(\Hom\) to
\[
 0\longrightarrow\Fp[[\mathbb Z_p]]
 \xrightarrow{\gamma-1}\Fp[[\mathbb Z_p]]
 \longrightarrow\Fp\longrightarrow0.
\]
This finite resolution compares continuously with the completed bar
resolution.  To check the assertion for the possibly noncompact
target \(M\), reduce modulo an invariant open subspace: a continuous
map from a profinite set then has finite image and extends linearly
to the free compact \(\Fp\)-module on that set.  Taking the inverse
limit over the open subspaces recovers \(M\) by completeness.
The same argument applies to
\(\operatorname{Map}_{\mathrm{cts}}(T,M)\), or directly to bar
cochains with parameter \(T\).  The finite analytic resolutions used
below have the same comparison property; see
Lemma~\ref{lem:finite-resolution} in
Section~\ref{sec:finite-resolutions}.

\subsection{Kummer annuli and the affine line}

Choose a compatible system of \(p\)-power roots of unity in \(C\).
It determines \(\epsilon\in C^\flat\) and a generator \(\gamma\) of
the Kummer group \(\mathbb Z_p(1)\).  Put
\(\theta=|\epsilon-1|\), so \(0<\theta<1\).
Let \(0<r\leq b\) belong to the value group of \(C\), and let
\(A[r,b]\) be the untilted closed annulus with coordinate \(z\).
We use the perfectoid coordinate-root construction and tilting
equivalence of \cite[Propositions~5.20 and~6.17]{Scholze12}.
After adjoining all \(p\)-power roots of \(z\), the perfectoid
Kummer cover has tilted function algebra
\begin{equation}
 \mathcal A[r,b]
 =
 \left\{\sum_{u\in\mathbb Z[1/p]}c_u z^u:
 \bigl(|c_u|\max(r^u,b^u)\bigr)_u\text{ is }c_0\right\}.
 \label{eq:analytic-annulus-algebra}
\end{equation}
The fractional monomials in this formula are coordinates on the
tilted cover.  The action is
\(\gamma(z^u)=\epsilon^u z^u\).
With parameter \(T\), replace \(c_u\in C^\flat\) by
\(c_u\in R_T:=\operatorname{Map}_{\mathrm{cts}}(T,C^\flat)\), using
the supremum norm and the uniform \(c_0\)-condition.

The cover and all its \v Cech levels are affinoid perfectoid.
Their higher structure-sheaf cohomology vanishes by
\cite[Proposition~8.8]{ScholzeDiamonds17}.
Kummer descent and \eqref{eq:analytic-cyclic} therefore identify the
cohomology of the annulus, including parameter \(T\), with the
cohomology of
\begin{equation}
 [\,\mathcal A[r,b](T)
          \xrightarrow{\gamma-1}\mathcal A[r,b](T)\,].
 \label{eq:analytic-annulus-complex}
\end{equation}

\begin{lemma}
\label{lem:annulus-restriction}
The degree-zero cohomology of
\eqref{eq:analytic-annulus-complex} is \(R_T\).
Suppose \(r\leq r'\leq b'\leq b\) and
\begin{equation}
 r'\geq r/\theta,\qquad b'\leq b\theta.
 \label{eq:analytic-radius-gap}
\end{equation}
After restriction to \(A[r',b']\), every degree-one class represented
by a Laurent series with zero constant coefficient vanishes.  Division by
\(\epsilon^u-1\), for \(u\neq0\), has operator norm at most one
between the indicated weighted coefficient spaces.  The remaining
degree-one class is the Kummer class of \(z\), with coefficients in
\(R_T\).
\end{lemma}

\begin{proof}
For \(u\neq0\), write \(u=p^v a\), where \(a\) is an integer prime to
\(p\).  Then
\begin{equation}
 |\epsilon^u-1|=\theta^{p^v},
 \qquad p^v\leq |u|_{\mathbb R}.
 \label{eq:analytic-epsilon}
\end{equation}
In particular, multiplication by \(\epsilon^u-1\) has zero kernel.
The invariants are therefore exactly the coefficient at \(u=0\).

For \(u>0\), the ratio of the target outer weight to the source
outer weight, after division, is bounded by
\[
 (b'/b)^u\theta^{-p^v}
 \leq\theta^{u-p^v}\leq1.
\]
For \(u<0\), the same calculation using the inner weights gives
\[
 (r'/r)^u\theta^{-p^v}
 \leq\theta^{|u|-p^v}\leq1.
\]
These inequalities prove both convergence and the asserted norm
bound, even for exponents with arbitrarily large \(p\)-power
denominators.  They preserve uniform \(c_0\)-tails with parameter
\(T\).  The constant summand of the two-term complex has zero
differential.  Its degree-one generator is the reduction modulo
\(p\) of the Kummer torsor of the coordinate, mapped into
\(\mathcal O^\flat\).
\end{proof}

The last assertion concerns the image of restriction; a fixed
closed annulus may have additional nonconstant degree-one classes.
This distinction will matter in the support calculation.

\begin{lemma}
\label{lem:affine-line-flat}
For every compact profinite \(T\),
\begin{equation}
 H^i_v\bigl(T\times(\mathbb A^1_C)^\diamond,\mathcal O^\flat\bigr)
 =
 \begin{cases}
 R_T,&i=0,\\
 0,&i>0.
 \end{cases}
 \label{eq:analytic-affine-line}
\end{equation}
The identification in degree zero is induced by constants.
\end{lemma}

\begin{proof}
First consider a closed disc \(D\).  Pullback to its perfectoid
power-map cover is injective on sections, by the sheaf property.
The pullback of a section is invariant under the root-of-unity
automorphisms.  On the punctured disc, the fractional monomial
description and \eqref{eq:analytic-epsilon} force all nonconstant
coefficients to vanish.  The same is then true on the whole
perfectoid disc.  Thus \(H^0(D^\diamond,\mathcal O^\flat)\) consists
of constants.  With parameter \(T\), these constants are precisely
\(R_T\).  This argument uses the power-map cover only for sections;
it does not treat the branched cover of a disc as a Kummer torsor.

Fix a disc \(D=\{|z|\leq R\}\).  Choose \(e\in C\) with
\begin{equation}
 |e|>R,\qquad R/|e|<|p|^{p/(p-1)}.
 \label{eq:analytic-distant-center}
\end{equation}
There are nested annuli about \(e\), both containing \(D\), whose
radii satisfy \eqref{eq:analytic-radius-gap}.  They are contained
in some larger disc \(D'\) about zero.  Restriction from \(D'\) to
\(D\) factors through these two annuli.  Since an annulus has
cohomology only in degrees \(0,1\), the factorization kills degrees
greater than one.  In degree one,
Lemma~\ref{lem:annulus-restriction} leaves only multiples of the
Kummer class of \(z-e\).

That remaining class vanishes on \(D\).  Indeed, choose a \(p\)-th
root of \(-e\) in \(C\).  The binomial series for
\((1-z/e)^{1/p}\) converges under
\eqref{eq:analytic-distant-center}: its \(n\)-th coefficient has
\(p\)-adic valuation \(-n-v_p(n!)\), so its radius of convergence is
\(|p|^{p/(p-1)}\).  It supplies a first \(p\)-th root of \(z-e\) on
\(D\).  This trivializes the mod-\(p\) Kummer torsor, hence also its
image in \(\mathcal O^\flat\)-cohomology.  The root and the annular
radius choices are independent of \(T\).

Now exhaust \(\mathbb A^1_C\) by a countable increasing sequence of
closed discs.  Derived global sections on the union are the derived
inverse limit of those on the discs.  In degree zero the restriction
system is the constant system \(R_T\).  In every positive degree it
is pro-zero by the preceding factorization.  The derived inverse
limit of a countable system has dimension at most one, so its
spectral sequence gives \eqref{eq:analytic-affine-line}.
\end{proof}

\subsection{Support on the locally profinite field}

We next calculate the fiber of restriction from the affine line to
the complement of \(\underline F\), where \(F/\Qp\) is unramified
quadratic.  All complements in this subsection are taken on
geometric fibers.

For a center \(c\in C\) and radius \(r>0\), write
\[
 E(c,r)=\{|z-c|\geq r\}\subset\mathbb A^1_C.
\]
It is the increasing union of closed annuli with inner radius \(r\)
and outer radius tending to infinity.  The tilted algebras of their
Kummer covers have dense restriction maps: fractional Laurent
polynomials are dense on every annulus and belong to every larger
annulus algebra.  Lemma~\ref{lem:analytic-roos}, in each bar degree,
therefore shows that the cohomology of \(T\times E(c,r)^\diamond\)
is computed by
\begin{equation}
 [\,\mathcal M(c,r;T)\xrightarrow{\gamma-1}
                  \mathcal M(c,r;T)\,],
 \label{eq:analytic-exterior-complex}
\end{equation}
where \(\mathcal M(c,r;T)\) is the Fr\'echet space of fractional
Laurent series converging on all these annuli, with the uniform
parameter convention.  Set
\[
 Q(c,r;T)=
 \coker\bigl(\gamma-1:\mathcal M(c,r;T)\to\mathcal M(c,r;T)\bigr).
\]
This is an ordinary algebraic cokernel; no assertion that the image
is closed is needed.

\begin{samepage}
\begin{lemma}
\label{lem:single-hole-residue}
The support complex
\[
 \operatorname{fib}\left(
 R\Gamma_v(T\times\mathbb A^{1,\diamond}_C,\mathcal O^\flat)
 \longrightarrow
 R\Gamma_v(T\times E(c,r)^\diamond,\mathcal O^\flat)\right)
\]
has cohomology only in degree \(2\), where it is \(Q(c,r;T)\).
Constant-coefficient extraction induces a surjection
\begin{equation}
 \operatorname{res}_{c,r}:Q(c,r;T)\longrightarrow R_T.
 \label{eq:analytic-residue}
\end{equation}
Its kernel maps to zero under restriction
\(Q(c,r;T)\to Q(c,r';T)\) whenever \(r'\geq r/\theta\).
Residues are preserved by enlarging a disc, by changing its center
inside the enlarged disc, and by scalar coordinate changes.
\end{lemma}
\end{samepage}

\begin{proof}
The invariants in \eqref{eq:analytic-exterior-complex} are \(R_T\).
They agree with the affine-line constants.  The support triangle
and Lemma~\ref{lem:affine-line-flat} therefore identify the sole
support group with \(Q(c,r;T)\) in degree \(2\).
The coefficient of \((z-c)^0\) vanishes on the image of
\(\gamma-1\); extraction is onto, with a lift of \(1\) given by the
Kummer class of \(z-c\).

If \(f\in\mathcal M(c,r;T)\) has constant coefficient zero, form
\begin{equation}
 h=\sum_{u\neq0}
       \frac{f_u}{\epsilon^u-1}(z-c)^u.
 \label{eq:analytic-exterior-primitive}
\end{equation}
For any target outer radius \(b\geq r'\), the negative exponents
are controlled by the inner-radius gap \(r'\geq r/\theta\);
the positive exponents are controlled by using the source outer
radius \(b/\theta\).  The calculation in
Lemma~\ref{lem:annulus-restriction} gives
\begin{equation}
 \|h\|_{[r',b]}\leq \|f\|_{[r,b/\theta]}.
 \label{eq:analytic-exterior-bound}
\end{equation}
The source is available for every \(b\), so \(h\) converges on the
whole enlarged exterior, with continuous parameter \(T\).  It
satisfies \((\gamma-1)h=f\) there.  This proves the uniform
vanishing of the residue kernel.

For completeness, the independence of the center is an assertion
about residues, and can be checked after making a further
enlargement.  If the centers are \(c,c'\), then sufficiently far
out
\[
 \frac{z-c'}{z-c}=1+\frac{c-c'}{z-c}
\]
has a first \(p\)-th root by the binomial estimate used in
\eqref{eq:analytic-distant-center}.  The two mod-\(p\) Kummer classes
then coincide.  Enlargement preserves the constant-coefficient
residue, and its kernels are killed by
\eqref{eq:analytic-exterior-bound}; hence the residue
identifications agree already before the further enlargement.
Multiplying a coordinate by a constant also preserves its Kummer
class, since a constant has a \(p\)-th root in \(C\).  The normal
Kummer generator is thus preserved by translations and by
multiplication by \(F^\times\).
\end{proof}

Choose a radius \(b_0\) in the value group with \(1<b_0<p\).
For \(n\geq0\), choose centers \(c_{n,\alpha}\in F\) for the cosets
\(\alpha\in F/p^n\mathcal O_F\), and put
\[
 r_n=|p|^n b_0.
\]
The discs of radius \(r_n\) about these centers are disjoint:
distinct centers have distance at least
\(|p|^{n-1}>r_n\).  Let \(E_n\) be their common exterior, including
the annular boundaries.  The exteriors are increasing with \(n\),
independently of choices of representatives, and
\begin{equation}
 (\mathbb A^1_C\setminus\underline F)^\diamond
      =\bigcup_{n\geq0} E_n^\diamond.
 \label{eq:analytic-tube-exhaustion}
\end{equation}
Here equality is also valid on affinoid perfectoid tests.
Write the untilt of such a test as \(S^\sharp=\Spa(A,A^+)\),
and write \(z\in A\) for its coordinate.  Give the uniform Banach
\(C\)-algebra \(A\) its spectral norm.  We use the Berkovich
spectrum \(\mathcal M(A)\) of bounded multiplicative seminorms
extending the norm of \(C\), with its seminorm topology.
It is compact Hausdorff: its defining relations cut out a closed
subset of \(\prod_{f\in A}[0,\|f\|]\).
Choose \(M>\max\{\|z\|,1\}\) and put
\(K=\{a\in F:|a|\leq M\}\), a compact subset of \(F\).
Points outside \(K\) cannot minimize the distance from \(z\).
For \(m\in\mathcal M(A)\), set
\[
 d(m)=\min_{a\in K}|z-a|_m.
\]
The ultrametric inequality gives, uniformly in \(m\),
\[
 \bigl||z-a|_m-|z-b|_m\bigr|\leq |a-b|.
\]
A finite \(\varepsilon\)-net of \(K\) therefore gives a
continuous finite minimum which approximates \(d\) uniformly
within \(\varepsilon\).  Thus \(d\) is continuous, without any
metrizability assumption on \(\mathcal M(A)\).
If \(d(m)=0\), then \(z=a\) in the completed residue field at
\(m\) for some \(a\in K\).  Passing to a completed algebraic
closure produces a geometric perfectoid fiber whose coordinate
lies in \(F\), contrary to fiberwise avoidance.  Compactness now
gives \(\delta=\min_m d(m)>0\).
Choose \(n\) with \(r_n<\delta\) strictly.  Every adic point
has a rank-one coarsening in \(\mathcal M(A)\).  The inequalities
\(|z-c_{n,\alpha}|>r_n\) hold at that coarsening and hence at
the original valuation: an inequality \(\leq\) cannot become
\(>\) on passing to a quotient of its ordered value group.
Thus the entire test, including its higher-rank points, factors
through \(E_n\).
Conversely, compatible membership in the shrinking tubes produces
compatible, locally constant cosets modulo \(p^n\mathcal O_F\).
Their limit is a continuous \(F\)-section, and the tube radii tending
to zero identify it with the given coordinate.  This proves
\eqref{eq:analytic-tube-exhaustion} with the stated fiberwise
meaning.

\begin{definition}
\label{def:locally-finite-distributions}
For a complete characteristic-\(p\) coefficient group \(R_0\), define
the locally finite distribution module on \(F\) by
\begin{equation}
 \mathcal D_{\mathrm{lf}}(F,R_0)
 =
 \varprojlim_n
 \prod_{\alpha\in F/p^n\mathcal O_F} R_0,
 \label{eq:analytic-distributions}
\end{equation}
where the transition sums over the finitely many children of each
coset.  Equivalently, take the product over \(F/\mathcal O_F\) of
the distribution modules on those compact open cosets.
Every finite-partition module and every product in
\eqref{eq:analytic-distributions} carries its product topology.
\end{definition}

\begin{proposition}
\label{prop:field-support}
For every compact profinite \(T\), the complex
\[
 \operatorname{fib}\left(
 R\Gamma_v(T\times\mathbb A^{1,\diamond}_C,\mathcal O^\flat)
 \longrightarrow
 R\Gamma_v(T\times
       (\mathbb A^1_C\setminus\underline F)^\diamond,
       \mathcal O^\flat)\right)
\]
has cohomology only in degree \(2\), where it is naturally
\(\mathcal D_{\mathrm{lf}}(F,R_T)\).
This identification is equivariant for translations and scalar
multiplication by \(F^\times\), with its natural action on
distributions.
\end{proposition}

\begin{proof}
Choose \(b_+\) in the value group with \(b_0<b_+<p\).  At tube
level \(n\), use the larger closed discs of radius
\(r_n^+=|p|^n b_+\) about all the \(c_{n,\alpha}\).
These discs and \(E_n\) form an open cover of the affine line in
the adic topology; the closed discs and annular boundaries here
are rational open subspaces.  The strict gap \(r_n<r_n^+\)
includes all boundary points.  In each bounded region only
finitely many discs occur.
Denote these pairwise disjoint larger discs by \(D_\alpha\), and
put \(V_n=\coprod_\alpha D_\alpha\).
For this paragraph abbreviate
\(R\Gamma_v(T\times Y^\diamond,\mathcal O^\flat)\) by
\(R\Gamma(Y)\).  The two-open cover
\(\mathbb A^1_C=E_n\cup V_n\) gives
\[
 R\Gamma(\mathbb A^1_C)\simeq
 \operatorname{fib}\left(
 R\Gamma(E_n)\oplus\prod_\alpha R\Gamma(D_\alpha)
 \longrightarrow
 \prod_\alpha R\Gamma(D_\alpha\cap E_n)\right).
\]
Taking the fiber of restriction to \(R\Gamma(E_n)\) yields
\[
 \prod_\alpha\operatorname{fib}\left(
 R\Gamma(D_\alpha)\longrightarrow
 R\Gamma(D_\alpha\cap E_n)\right).
\]
Here products are ordinary products of complexes of
\(\Fp\)-vector spaces and are exact.  For a single hole, the
two-open cover by \(D_\alpha\) and its full exterior identifies
the corresponding factor with the support complex of
Lemma~\ref{lem:single-hole-residue}.  Thus the sole cohomology group
at this stage is
\begin{equation}
 Q_n(T)=
 \prod_{\alpha\in F/p^n\mathcal O_F}
       Q(c_{n,\alpha},r_n;T)
 \quad\text{in degree }2.
 \label{eq:analytic-level-support}
\end{equation}
This argument uses an actual open cover and ordinary sheaf
cohomology; in particular the product at infinity has no imposed
uniform norm bound.

The inclusions \(E_n\subset E_{n+1}\) induce an inverse system on
the support complexes.  A parent receives the sum of the support
maps from its finitely many children.  By
Lemma~\ref{lem:single-hole-residue}, the residue of this map is
literally the sum of their residues.  We obtain short exact
sequences of inverse systems
\begin{equation}
 0\longrightarrow K_n(T)\longrightarrow Q_n(T)
 \longrightarrow
 D_n(T):=\prod_{\alpha\in F/p^n\mathcal O_F}R_T
 \longrightarrow0.
 \label{eq:analytic-residue-system}
\end{equation}
It is not necessary to choose Kummer lifts commuting with
refinement: their residues, rather than their lifts, are
canonical.

Choose once and for all an integer \(s\geq1\) such that
\(p^s\geq\theta^{-1}\).  A level-\(n+s\) child has radius
\(r_{n+s}\), and its parent's radius is \(r_n=p^s r_{n+s}\).
The child center lies strictly inside the parent disc.  The parent
exterior is therefore the same exterior when recentered at that
child.  In the child-centered Kummer complex for this same parent
exterior, a residue-kernel representative has a primitive given
by \eqref{eq:analytic-exterior-bound}.  It is therefore an actual
boundary in the ordinary algebraic cokernel.  Its vanishing is
intrinsic to the parent exterior and does not require the
parent's original centered Kummer cover to agree with the
child's.  Consequently the following map is zero:
\[
 K_{n+s}(T)\longrightarrow K_n(T).
\]
This same \(s\) works for all \(n\), all centers, all
parents in the product, and all \(T\).  Thus the kernel system is
uniformly pro-zero.

The transition \(D_{n+1}(T)\to D_n(T)\) is split surjective:
choose one child of each parent, put the parent coefficient in
that child, and put zero in its other children.  This defines a
continuous section for the product topologies.  Hence its derived
inverse limit is its ordinary inverse limit
\(\mathcal D_{\mathrm{lf}}(F,R_T)\), whereas the derived inverse
limit of \(K_n(T)\) is zero.  Formula
\eqref{eq:analytic-tube-exhaustion}, the support triangles, and
exactness of limits in the derived category now prove the
assertion.  The one-hole annular exhaustions involved in this
argument have no additional derived-limit contribution by
Lemma~\ref{lem:analytic-roos} and
\eqref{eq:analytic-exterior-complex}.

Translations and scalar multiplication send tube systems to
cofinal tube systems.  Their action on the residue targets is the
one proved in Lemma~\ref{lem:single-hole-residue}; consequently the
limit action is the natural action on distributions.
\end{proof}

\subsection{The translation quotient and its orientation}

The shifted term has a derived Steinberg antecedent. Heyer's published
calculation treats $\mathrm{GL}_2(\mathbb Q_p)$ with algebraically closed
characteristic-$p$ coefficients \citep[Corollary~5.3.4]{Heyer23}; his later
geometrical lemma treats finite extensions of $\mathbb Q_p$
\citep[Lemma~4.1.6 and Corollary~4.3.10]{HeyerGeo24}.
The comparison of smooth and solid coefficients in the coheight-one setting
is made in \citet[Sections~3.5--3.6]{BMR26}.
We give the rank-two continuous group-cohomology calculation directly,
retaining the topologies in Definition~\ref{def:locally-finite-distributions}
and the action on the orientation line.

\begin{lemma}
\label{lem:translation-cohomology}
Let \(R_0=C^\flat\) or \(R_T\), with trivial additive \(F\)-action.
Then
\begin{equation}
 H^i_{\mathrm{cts}}(F,R_0)=
 \begin{cases}R_0,&i=0,\\0,&i>0,\end{cases}
 \label{eq:analytic-constant-translations}
\end{equation}
and
\begin{equation}
 H^i_{\mathrm{cts}}(F,\mathcal D_{\mathrm{lf}}(F,R_0))=
 \begin{cases}R_0,&i=2,\\0,&i\neq2.\end{cases}
 \label{eq:analytic-distribution-translations}
\end{equation}
In \eqref{eq:analytic-distribution-translations},
multiplication by \(u\in\mathcal O_F^\times\) acts by
\(N_{F/\Qp}(u)^{-1}\bmod p\), under the convention of pushforward
on distributions and conjugation on cochains.
\end{lemma}

\begin{proof}
Write \(F=\bigcup_{m\geq0}p^{-m}\mathcal O_F\).
Each lattice is isomorphic to \(\mathbb Z_p^2\).
The completed group-ring resolution in two coordinates computes
its continuous cohomology with trivial \(R_0\)-coefficients as
\(\bigwedge^\ast(R_0^2)\).  Restriction from one lattice to its
\(p\)-multiple is zero in every positive degree: on degree one it
is multiplication by \(p=0\), and the higher-degree maps are its
exterior powers.

At the level of bar cochains, restriction between these compact
open lattices is surjective, by extending a continuous function
by zero on the clopen complement.  Consequently continuous
cochains on \(F\) are their derived inverse limit, including
parameter \(T\).  The degree-zero cohomology system is constant
and the positive systems are pro-zero.  This proves
\eqref{eq:analytic-constant-translations}.

Set \(\Lambda=\mathcal O_F\), considered additively.  By the
product over \(\Lambda\)-cosets,
\(\mathcal D_{\mathrm{lf}}(F,R_0)\) is the continuous coinduction
from \(\Lambda\) of \(\mathcal D(\Lambda,R_0)\).
The quotient \(F/\Lambda\) is discrete; any set of representatives
is therefore a continuous section.  The usual bar comparison for
Shapiro's lemma, and its simplicial contraction, are continuous
with this section and the product topology.  It reduces
\eqref{eq:analytic-distribution-translations} to
\(\Lambda\)-cohomology.

Choose a \(\mathbb Z_p\)-basis of \(\Lambda\).
At finite quotient level the group-ring coordinate maps
\(\gamma_i-1\) identify its distribution module with
\[
 R_0[X_1,X_2]/(X_1^{p^n},X_2^{p^n}).
\]
The transition is the quotient map, because it sums coefficients
over fibers of the finite group quotient.  Inverse limit gives
\[
 \mathcal D(\Lambda,R_0)=R_0[[X_1,X_2]]
\]
with the product topology on its coefficients.  This is the
formal power-series module, without restrictions on its
coefficient sequence.

The two-coordinate completed resolution now gives the
cohomological Koszul complex
\begin{equation}
 0\longrightarrow M
 \xrightarrow{(X_1,X_2)}
 M^2
 \xrightarrow{(-X_2,X_1)}
 M\longrightarrow0,\qquad M=R_0[[X_1,X_2]].
 \label{eq:analytic-koszul}
\end{equation}
Multiplication by \(X_1\) is injective, and multiplication by
\(X_2\) is injective after reduction modulo \(X_1\).
Thus the only cohomology is
\(M/(X_1,X_2)=R_0\) in degree \(2\).
Coefficient deletion, insertion, and shifts give continuous maps
in this argument, so it also proves the statement with \(R_T\)
and with further compact profinite parameters.

Finally, a lattice automorphism with matrix \(A\in
\operatorname{GL}_2(\mathbb Z_p)\) induces on the linear terms of
the completed group coordinates the matrix \(A\bmod p\).
Its action on the top exterior generator of the resolution is
\(\det A\bmod p\).  Applying \(\Hom\), with the contravariant
cochain convention, gives the inverse determinant on the
augmentation quotient.  For multiplication by \(u\), this
determinant is \(N_{F/\Qp}(u)\).  The construction is independent of
the chosen lattice basis and agrees under Shapiro's comparison.
\end{proof}

\begin{proposition}
\label{prop:geometric-cohomology}
For every compact profinite \(T\), the geometric period domain has
cohomology
\begin{equation}
 H^i_v(T\times(N^\ast)_{C^\flat},\mathcal O^\flat)
 =
 \begin{cases}
 R_T,&i=0,\\
 R_T\otimes_{C^\flat}\ell_C,&i=3,\\
 0,&i\neq0,3,
 \end{cases}
 \label{eq:analytic-geometric-rows}
\end{equation}
where \(\ell_C\) is a one-dimensional \(C^\flat\)-space.
All identifications are natural under continuous maps of \(T\)
and the frame-group actions.  On \(\ell_C\),
\(\mathcal O_F^\times\subset J\) acts through the nontrivial
character
\[
 \delta(u)=N_{F/\Qp}(u)\bmod3.
\]
The kernel-scalar action of \(\Zp^\times\) is trivial.
The line spaces in \eqref{eq:analytic-geometric-rows} have their
usual topology, as expressed by the continuous-parameter formula.
\end{proposition}

\begin{proof}
By Lemma~\ref{lem:negative-presentation},
\[
 (N^\ast)_{C^\flat}
 \simeq
 [\,(\mathbb A^1_C\setminus\underline F)^\diamond/
                  \underline F\,].
\]
The \v Cech spectral sequence of the translation quotient may be
computed by continuous \(F\)-cochains.  Indeed, in each bar degree
its locally profinite parameter \(F^r\) is a disjoint union of
compact open pieces.  The preceding analytic calculations apply
on every such piece, and ordinary products give their global
cochains.

Apply continuous \(F\)-cohomology to the support triangle of
Proposition~\ref{prop:field-support}.
Lemma~\ref{lem:affine-line-flat} and
\eqref{eq:analytic-constant-translations} put the affine-line
term in degree zero, equal to \(R_T\).
The support term is in degree \(2\) before group cohomology, and
\eqref{eq:analytic-distribution-translations} puts it in degree
\(4\) afterward.  The support triangle therefore puts the
additional class of the complement quotient in degree \(3\).
This proves the degrees and dimensions in
\eqref{eq:analytic-geometric-rows}.

Scalar multiplication preserves the normal Kummer residue.
The extra action is consequently the rank-two lattice orientation
of Lemma~\ref{lem:translation-cohomology}.
At \(p=3\) its inverse equals itself, giving \(\delta\).
The residue norm \(\mathbb F_9^\times\to\mathbb F_3^\times\) is
surjective, so this character is nontrivial.
A scalar \(a\in\Zp^\times\), regarded as an element of \(F^\times\),
has norm \(a^2\), which is \(1\bmod3\).  Thus the determinant
adjustments of the kernel frame in
Lemma~\ref{lem:negative-presentation} act trivially on the line.
Inversion of frame conventions and the other unramified embedding
give the same order-two character.

All coefficient contractions and all lattice calculations were
made with \(R_T\), with their supremum and product topologies.
Evaluation on \(T\) therefore identifies the results with
continuous functions to the displayed lines.  No equivariant
splitting of the complex between its two cohomology degrees is
asserted or needed.
\end{proof}

\subsection{Frobenius on actual fixed-stage complexes}

The geometric calculation has produced two $C^\flat$-lines with their
frame actions.  We now descend their coefficients to $k$, while also
identifying the arithmetic $H$-cochains of each finite marking stage
$B=L^U$.  Keeping $B$ fixed is essential: the quotient by $U$ can
contribute additional cohomology, whose finiteness is needed in
Section~\ref{sec:decompletion}.  The marking colimit will be taken only
after that finite-stage calculation.

Put \(q=|k|\).  Base Frobenius on \(C^\flat\) gives a map
\(\Phi\) on spaces and complexes defined over \(k\).  It is
defined before choosing the untilt presentation in
Lemma~\ref{lem:negative-presentation}.  In particular the chosen
untilt need not be fixed.

For a finite marking stage \(B=L^U\), set
\(X_B=\Spa(B^\flat)\), and let
\[
 \mathcal A_B
   =H^0_v(X_B\times_k\Spa(C^\flat),\mathcal O^\flat).
\]
Geometric base change here uses the v-sheaf associated with the
finite-field base.  Each field factor of \(B\) is a complete
equal-characteristic local field with finite residue field, hence
has the form \(k'((s))\).
After splitting its finite constants, its geometric base change
is the perfected punctured open unit disc.  This identification
can be checked directly on a characteristic-\(p\) perfectoid test
algebra over \(C^\flat\): a continuous map from
\(\widehat{k((s))^{\mathrm{perf}}}\) is specified by the image of
\(s\), an invertible topologically nilpotent element, together
with its unique \(p\)-power roots.  Conversely such an element
evaluates every completed fractional Laurent series, by
convergence of its tails.  These are precisely the points of the
perfected punctured open disc.  Thus \(\mathcal A_B\)
is a finite product of fractional Laurent Fr\'echet algebras on
\(0<|s|<1\), with constants in \(C^\flat\).
Affinoid perfectoid acyclicity and
Lemma~\ref{lem:analytic-roos} show that this geometric space, also
after multiplication by a compact profinite set, has no higher
structure-sheaf cohomology.

\begin{lemma}
\label{lem:frobenius-descent}
Let \(B\) be a finite product of finite separable extensions of
\(k((t))\).  For every compact profinite \(Q\) there is an exact
sequence
\begin{equation}
 0\longrightarrow
 \operatorname{Map}_{\mathrm{cts}}(Q,B^\flat)
 \longrightarrow
 \operatorname{Map}_{\mathrm{cts}}(Q,\mathcal A_B)
 \xrightarrow{\Phi-1}
 \operatorname{Map}_{\mathrm{cts}}(Q,\mathcal A_B)
 \longrightarrow0.
 \label{eq:analytic-frobenius-exact}
\end{equation}
It is natural in \(B,Q\) and continuous actions over \(k\).
In particular, for a compact group \(G\) acting continuously on
\(B\), over \(k\) and preserving its valuations, the natural map
identifies actual cochain complexes by
\begin{equation}
 \Ccts(G,B^\flat;T)
 \simeq
 \operatorname{fib}\left(
 \Phi-1:\Ccts(G,\mathcal A_B;T)
               \longrightarrow\Ccts(G,\mathcal A_B;T)\right).
 \label{eq:analytic-frobenius-cochains}
\end{equation}
\end{lemma}

\begin{proof}
Begin with \(B=k((s))\).  Elements of \(\mathcal A_B\) are
fractional Laurent series
\(\sum_{u\in\mathbb Z[1/p]}c_u s^u\) converging on every smaller
closed annulus in \(0<|s|<1\).
Frobenius sends \(c_u\) to \(c_u^q\) and leaves \(s\) fixed.
Its fixed coefficients belong to \(k\).
The annular \(c_0\)-condition for finite-field coefficients says
exactly that the support is bounded below and has only finitely
many exponents below any fixed real bound.  These are the series
in \(\widehat{k((s))^{\mathrm{perf}}}=B^\flat\).
Moreover, the topology induced by the annular seminorms is its
valuation topology: for nonzero finite-field coefficients the
norm of a series is determined by its least exponent.  This
identifies the kernel in \eqref{eq:analytic-frobenius-exact},
including its topology.

We prove surjectivity with the parameter \(Q\) present.
A continuous \(Q\)-family has continuous coefficient functions
\(c_u:Q\to C^\flat\) and uniformly small weighted tails in each
annulus seminorm.  For each \(c\in C^\flat\), there is \(b\in C^\flat\)
satisfying
\begin{equation}
 b^q-b=c.
 \label{eq:analytic-artin-schreier}
\end{equation}
If \(|c|<1\), the unique solution with \(|b|<1\)
is
\[
 b=-\sum_{j\geq0}c^{q^j},\qquad |b|=|c|.
\]
If \(|c|>1\), every solution has
\(|b|=|c|^{1/q}\); if \(|c|=1\), a solution has norm at most one.
In every case we can choose
\begin{equation}
 |b|\leq |c|.
 \label{eq:analytic-root-bound}
\end{equation}
The polynomial in \eqref{eq:analytic-artin-schreier} has derivative
\(-1\).  Near each value \(c_0\), a chosen root therefore extends
to a continuous local branch, obtained by adding the small
solution for \(c-c_0\).
For a continuous coefficient \(c_u:Q\to C^\flat\), compactness of
\(Q\) and a finite clopen refinement give a continuous choice
\(b_u\) with \eqref{eq:analytic-root-bound}.
The partitions may depend on \(u\).
The uniform bound preserves the weighted \(c_0\)-tails in every
annulus seminorm.  More explicitly, fix one such seminorm, write
its monomial weight as \(w(u)\), and fix \(q_0\in Q\) and
\(\varepsilon>0\).  Outside a finite set of exponents \(I\),
we have \(\sup_{q\in Q}|b_u(q)|w(u)<\varepsilon\).
The same bound holds for
\(|b_u(q)-b_u(q_0)|w(u)\) by ultrametricity.  For the finitely
many \(u\in I\), intersect the continuity neighborhoods of
the coefficient functions.  This proves continuity in the
chosen annular seminorm; applying it to each seminorm proves
continuity into \(\mathcal A_B\).  Thus the assembled series
is a continuous solution of
\eqref{eq:analytic-frobenius-exact}, although the partitions
chosen for its different coefficients need not agree.

For \(B=k'((s))\), put \(f=[k':k]\).  Geometric base change
has one factor for each \(k\)-embedding of \(k'\) into
\(C^\flat\).  Index these factors modulo \(f\), so that
\(\Phi(x)_i=F(x_{i-1})\), where \(F\) raises each Laurent
coefficient to its \(q\)-th power.  To solve
\(\Phi(x)-x=y\), first solve
\[
 (F^f-1)x_0
   =y_0+\sum_{j=1}^{f-1}F^{f-j}y_j,
\]
using the preceding argument with \(q\) replaced by \(q^f\),
and then put \(x_i=F(x_{i-1})-y_i\) for
\(1\leq i<f\).  These operations preserve the Fr\'echet
algebra and all continuous parameter families, since
\[
 \|F^j a\|_{[r,b]}
   =\|a\|_{[r^{1/q^j},b^{1/q^j}]}^{q^j}.
\]
All transformed radii remain inside the open unit interval.
In the kernel, \(F^f x_0=x_0\) and \(x_i=F^i x_0\),
which recovers the original finite constant field.
Every nonzero finite-field coefficient has norm one, so the
annular norms on this kernel depend only on the least exponent
and induce exactly the valuation topology of \(B^\flat\).
Finite cycles and finite products preserve this identification.
The resulting kernel inclusion is intrinsic to \(B\), and
is independent of the uniformizer and of the chosen decomposition
of the constants.

The auxiliary root choices prove surjectivity; they are not
claimed to be equivariant.  The kernel inclusion and
\(\Phi-1\) themselves are canonical and natural.

Finally apply \eqref{eq:analytic-frobenius-exact} with
\(Q=T\times G^r\) in every bar degree.
The maps commute with the bar faces, since the action is over
\(k\) and commutes with base Frobenius.
Termwise exactness gives \eqref{eq:analytic-frobenius-cochains}.
\end{proof}

Let
\begin{equation}
 K_C(T)=R\Gamma_v(T\times(N^\ast)_{C^\flat},\mathcal O^\flat),
 \qquad
 K_k(T)=\operatorname{fib}(\Phi-1:K_C(T)\to K_C(T)).
 \label{eq:analytic-arithmetic-complex}
\end{equation}
On a one-dimensional geometric cohomology line, \(\Phi\) is a
bijective \(q\)-semilinear map.  Writing it as
\(ae\mapsto c a^q e\), one obtains a nonzero fixed basis by solving
\(c a^{q-1}=1\).  In that basis, its difference with the identity is
\eqref{eq:analytic-artin-schreier}.
The proof with continuous parameters in
Lemma~\ref{lem:frobenius-descent} applies to these lines.
Proposition~\ref{prop:geometric-cohomology} and the long exact
sequence of the fiber consequently give
\begin{equation}
 H^i K_k(T)=
 \begin{cases}
 \operatorname{Map}_{\mathrm{cts}}(T,k),&i=0,\\
 \operatorname{Map}_{\mathrm{cts}}(T,\ell),&i=3,\\
 0,&i\neq0,3,
 \end{cases}
 \label{eq:analytic-arithmetic-rows}
\end{equation}
where \(\ell\) is a one-dimensional \(k\)-space.
The first identification is induced by constants.
The character of \(\mathcal O_F^\times\) on \(\ell\) is \(\delta\),
and the kernel-scalar action is trivial.

The frame-group action on \(\ell\) is continuous for the discrete
topology on \(k\).  Indeed, apply
\eqref{eq:analytic-arithmetic-rows} to the universal compact
profinite family of frame changes.  The resulting functions take
values continuously in the finite fixed line.  In particular this
is a finite smooth \(J\)-module.

\subsection{Finite marking stages and their colimit}

The arithmetic complex $K_k(T)$ has the two discrete cohomology lines
in \eqref{eq:analytic-arithmetic-rows}.  At a fixed stage $B=L^U$,
Theorem~\ref{thm:two-towers} identifies the required $H$-quotient
with $[N^*/U]$.  For sufficiently small $U$, we first prove finiteness
and identify the cohomology with continuous profinite parameters.
Shrinking $U$ then removes its positive-degree cohomology and leaves
the two lines in the marking colimit.

\begin{theorem}
\label{thm:completed-calculation}
For a cofinal family of sufficiently small open normal subgroups
\(U\subset J\), put \(B=L^U\).
Then \(H^a_{\mathrm{cts}}(H,B^\flat)\) is finite for every \(a\),
and it is zero for \(a>7\).
For every compact profinite \(T\), evaluation induces the natural
identification
\begin{equation}
 H^a\Ccts(H,B^\flat;T)
 =
 \operatorname{Map}_{\mathrm{cts}}
       \bigl(T,H^a_{\mathrm{cts}}(H,B^\flat)\bigr),
 \label{eq:analytic-finite-stage-parameters}
\end{equation}
where the finite group on the right is discrete.
Moreover,
\begin{equation}
 (B^\flat)^H=k,\qquad L^H=k,
 \label{eq:analytic-invariant-field}
\end{equation}
and
\begin{equation}
 \colim_U H^a\Ccts(H,(L^U)^\flat;T)
 =
 \begin{cases}
 \operatorname{Map}_{\mathrm{cts}}(T,k),&a=0,\\
 \operatorname{Map}_{\mathrm{cts}}(T,\ell),&a=3,\\
 0,&a\neq0,3.
 \end{cases}
 \label{eq:analytic-completed-colimit}
\end{equation}
The colimit identifications respect restrictions, conjugations and
the \(J\)-action.  The residual \(P/H=\Zp^\times\)-action on both
lines is trivial; on the upper line
\(\mathcal O_F^\times\subset J\) acts by \(\delta\).
The statements are natural in \(T\).
\end{theorem}

\begin{proof}
\emph{The fixed-stage comparison.}
At a fixed stage \(B=L^U\), Theorem~\ref{thm:two-towers} gives
\begin{equation}
 [X_B/H]\simeq[N^\ast/U].
 \label{eq:analytic-fixed-stage-quotient}
\end{equation}
It is an identity of quotient v-stacks with their actual
restriction and frame-change maps.
After geometric base change, take derived global sections,
including the product with \(T\), and then the fiber of
\(\Phi-1\).
On the left,
Lemma~\ref{lem:frobenius-descent} in every \(H\)-bar degree gives
exactly \(\Ccts(H,B^\flat;T)\).
On the right, fibers commute with totalization, giving the
\(U\)-bar totalization of the complexes \(K_k(T\times U^r)\).
Thus there is a natural comparison
\begin{equation}
 \Ccts(H,B^\flat;T)
 \simeq
 \Tot\bigl([r]\longmapsto K_k(T\times U^r)\bigr).
 \label{eq:analytic-stage-comparison}
\end{equation}
Every operation here is at the fixed stage \(B\).
The comparison uses the pro-\'etale torsor in
Theorem~\ref{thm:two-towers}; it involves no interchange between
continuous cochains and completion of the entire marking union.

\smallskip\noindent\emph{Finiteness and profinite families.}
Choose \(U\) small enough to be torsion-free uniform and to fix the
finite lines in \eqref{eq:analytic-arithmetic-rows}.
Such \(U\) form a cofinal family; they may be chosen normal in
\(J\).
Their Lie dimension is \(4\).
Continuous cohomology with finite characteristic-\(p\)
coefficients is finite and vanishes above degree \(4\), by the
finite analytic resolution; see
Lemma~\ref{lem:finite-resolution}.
The hypercohomology spectral sequence of
\eqref{eq:analytic-stage-comparison} has only the rows
\begin{equation}
 E_2^{r,0}=H^r_{\mathrm{cts}}(U,k),\qquad
 E_2^{r,3}=H^r_{\mathrm{cts}}(U,\ell),
 \label{eq:analytic-stage-spectral-sequence}
\end{equation}
and \(0\leq r\leq4\).
It is first-quadrant with a finite possible range.
Therefore the abutment is finite and vanishes above degree \(7\).
No splitting between the two rows is needed for this conclusion.

We verify the family assertion in this spectral sequence.
For finite discrete coefficients \(V\), a continuous function
\(T\times U^r\to V\) factors through a finite clopen partition of
\(T\): cover the compact product by finitely many clopen
rectangles on which the function is constant, and refine the
partition of \(T\).
Thus cycles and boundaries in these row complexes may be
computed on finitely many \(T\)-pieces.
The same applies to primitives, and gives
\[
 H^r_{\mathrm{cts}}\!
       \left(U,\operatorname{Map}_{\mathrm{cts}}(T,V)\right)
 =
 \operatorname{Map}_{\mathrm{cts}}
       \left(T,H^r_{\mathrm{cts}}(U,V)\right)
\]
in the indicated parameter convention.
Together with \eqref{eq:analytic-arithmetic-rows}, this gives
\(E_2^{r,j}(T)=\operatorname{Map}_{\mathrm{cts}}
(T,E_2^{r,j}(*))\).
All these groups at a point are finite.  At every subsequent
page, naturality under point evaluations makes the differential
pointwise equal to the differential at a point: evaluation in
its target is injective by the preceding-page identification.
Continuous functions on \(T\) are exact on finite discrete
groups, by lifting on a finite clopen partition.  Taking kernels
and cokernels therefore proves the same identification on the
next page.  The fixed finite rectangle of possible bidegrees
gives a uniform limiting page and its natural finite abutment
filtration.
Here is also a direct verification for the abutment extensions.
Suppose a functorial extension has outer terms
\(\operatorname{Map}_{\mathrm{cts}}(T,A)\) and
\(\operatorname{Map}_{\mathrm{cts}}(T,B)\), with \(A,B\) finite,
and its middle functor commutes with finite clopen decompositions
of \(T\).  On a partition where the image in \(B\) is constant,
lift that constant using an element of the middle group at a
point.  Subtracting its constant pullback leaves a continuous
\(A\)-valued function.  Evaluation in the middle group is
therefore locally constant.  The same construction realizes
every locally constant function into the middle group at a
point, and evaluation is injective by the two outer terms.
All the filtration functors here commute with finite clopen
decompositions, since the original complexes do.
Induction through the finite filtration gives
\eqref{eq:analytic-finite-stage-parameters}.
This argument establishes the finite-stage parameter assertion
before any use of strictness in the uncompleted calculation.

\smallskip\noindent\emph{Invariants and the marking colimit.}
In total degree zero, the sole term of
\eqref{eq:analytic-stage-spectral-sequence} is \(H^0(U,k)=k\).
It is induced by the constant unit.  Hence
\((B^\flat)^H=k\).
The inclusions \(B\hookrightarrow B^\flat\) are injective and
\(H\)-equivariant, so \(B^H=k\) as well.
Every element of \(L\) lies in some stage and then in a
sufficiently large stage in the chosen cofinal family.
This proves \(L^H=k\) without any decompletion argument.

It remains to pass to the marking colimit.
For a finite smooth coefficient module \(V\), every
positive-degree continuous cohomology class of an open subgroup
becomes zero after restriction to a sufficiently small open
subgroup.  One may use normalized inhomogeneous cochains:
a continuous cocycle is zero at the identity tuple and has
discrete values, hence is zero on the power of a sufficiently
small open subgroup.
Degree-zero invariants have colimit \(V\).
The spectral sequences
\eqref{eq:analytic-stage-spectral-sequence} are uniformly bounded,
and filtered colimits of abelian groups are exact.
Their colimit therefore leaves only \(k\) in degree zero and
\(\ell\) in degree three, proving
\eqref{eq:analytic-completed-colimit} for a point.

For general \(T\), use
\eqref{eq:analytic-finite-stage-parameters}.
Continuous maps to a discrete group have finite image, so
\(\operatorname{Map}_{\mathrm{cts}}(T,-)\) commutes with a
filtered colimit of discrete groups: finitely many values and
finitely many equalities are realized at one common stage.
This gives \eqref{eq:analytic-completed-colimit} with all \(T\).

Finally, \eqref{eq:analytic-fixed-stage-quotient} and
\eqref{eq:analytic-stage-comparison} are natural for all frame
changes.  Their colimit action is therefore the action on the two
lines in \eqref{eq:analytic-arithmetic-rows}.
By Proposition~\ref{prop:geometric-cohomology}, the residual
determinant acts through scalar multiplication by a
\(\Zp^\times\)-unit on the negative space, with trivial action on
both lines at \(p=3\).  The \(J\)-action restricts on
\(\mathcal O_F^\times\) to \(\delta\) on \(\ell\), as claimed.
\end{proof}

The output of this section concerns the completed perfections of
individual finite marking stages.  The next section compares
their cochain colimit with the root-stage cochains specified in
Definition~\ref{def:stage-cochains}.
